\documentclass[dvipdfmx]{jlreq}
\usepackage{graphicx}
\usepackage{float}

% \title{粘土による直感的操作と生成AIを併用したデザイン支援手法}
\title{粘土による直感的操作と生成AIを併用したデザイン探索手法}
\author{梅木敦成}
\date{\today}

\begin{document}
\maketitle


\section{はじめに}
\begin{itemize}
  \item ものづくり、デザインでの創作が背景
  \item 多くの非専門家は、「なんとなくこういうものが作りたい」という初期イメージ(種)を持ってはいるものの、それは輪郭が不明瞭で曖昧な状態であることが多い
  \item 非専門家にとって頭の中のイメージを具体化, 具現化にするのは簡単ではない
  \item 従来はスケッチやCADツールが用いられてきたが、これらは高い操作スキルや描画能力を必要とするため、非専門家には難しいプロセスとなっていた。
  \item しかし、近年 Text-to-Image 生成AIが、この「描画スキル」を担ってくれることが期待されている。。 
  \item つまり、言語（プロンプト）インターフェースによって、描画や造形のプロセスを生成 AI に任せられるようになると期待されている。
  \item 
  \item 【課題：パターンの空間と言葉の空間の乖離】
  \item しかし、言語による指示は、生成AIを利用して自身のイメージを具現化しようとする利用者（ユーザー）にとって、新たな課題をもたらしている。
  \item 多くのユーザーはプロンプトエンジニアリング（試行錯誤）によって解決を図っている。
  \item 言葉は「雰囲気（スタイル）」や「概念」を伝えることには適しているが、「ここをこのように曲げたい」「この部分のボリュームを増やしたい」といった具体的な形状（ジオメトリ）や空間的な特徴を正確に記述することには本質的に不向きである。
  \item その結果、ユーザーは言語化できない形状へのこだわりを妥協するか、あるいは偶然生成される出力に依存せざるを得ない状況にある。
  \item 
  \item ＊研究の目指すこと＊ (社会的意義なども含め)
  \item 本研究は、非専門家ユーザーの「曖昧で多義的な初期イメージ」を、システムとの対話を通じて「言語化・具体化」し、最終的に「明確な要求仕様」へと昇華させることを目標とする。
  \item 本研究の社会的意義
  \begin{itemize}
    \item 「創造性の民主化」の推進：高度な造形スキルを持たない人々が、自身の内面にあるイメージを正確に外部化できるよう支援することで、個人の創造的な表現活動を促進する。
    \item 「身体性とデジタル技術の融合」：AIによる自動化が進む現代において、人間が本来持っている「手で考える（触覚的思考）」能力を排除するのではなく、AIによる解釈支援と組み合わせることで、人間の身体的直感とAIの表現力を共存させる新たな創作プロセスを提示する。
  \end{itemize}
  \item 
  % \item *問題点＊
  % \item 非専門家は、「なんとなくこういうものが作りたい」という初期イメージ（種）は持っているが、それは不明瞭な状態である。
  \item 
  \item ＊研究のアイディアをごく簡潔に＊
  \item そこで本研究では、生成AIにおける「言語化の壁」を突破するために、言語化しにくい形状特徴を直接入力できる（粘土などの）物理的なインターフェースと、AIによる解釈支援を組み合わせた手法を提案する。
\end{itemize}



% \section{提案手法}
\section{生成 AI を用いた物品のデザイン}
\subsection{デザインにおける生成 AI の利用}
% \begin{itemize}
%   \item ＊問題の分析＊
%   \item ＊造形にこだわる必要はない＊ (造形に特化して十分に進捗しているのであれば良いが，アイディア先行の状態なので，限定するのは逆効果)
%   \item 本来、制作過程で試行錯誤しながら頭の中の不明瞭なイメージを明確化（Solidify）していく必要があるが、非専門家は専門的な語彙に乏しいため、AIに対して的確な指示が出せず、イメージを固めるための対話が成立しない。
%   \item 結果として、多義的な言葉（「かっこいい」等）に頼らざるを得ず、「自分の意図（メンタルモデル）と出力の乖離」を埋められないまま、偶然の正解を待つことになる。（＝意図通りのデザインが出力されない、効率が悪い）
%   \item さらに、言葉は「雰囲気（スタイル）」を伝えるのは得意だが、「具体的な形状（ジオメトリ）」を伝えるのには本質的に不向きである。（例：「ここをこう曲げて」はテキストで書きにくい）。
%   \item 既存手法（Text-to-Image）では、この「形状へのこだわり」を反映させる手段が欠如している。
%   \item 既存の Image-and-Text-to-Image でオンラインの画像で検索して補完する手法もあるが、ユーザーがイメージする形状特徴を持つ画像が必ずしも存在するとは限らない。
% \end{itemize}
\begin{itemize}
  \item \textbf{デザインプロセスにおける「探索」の重要性}
  \begin{itemize}
    \item デザインとは、完成された設計図を最初から出力する作業ではなく、試行錯誤を通じて頭の中の不明瞭なイメージを徐々に明確化（Solidify）していく「探索的プロセス」である。
    \item 特に初期段階においては、ユーザー自身も最終的な正解を持っておらず、AIからのフィードバックを見ながら自身の意図（メンタルモデル）を更新していく対話的なプロセスが重要となる。
  \end{itemize}

  \item \textbf{言語インターフェースの限界（情報の非対称性）}
  \begin{itemize}
    \item 現在主流のText-to-Image手法は、自然言語を入力とする。言葉は「雰囲気（スタイル）」や「概念（コンセプト）」といった抽象的な情報を伝えることには長けている。
    \item しかし、デザインにおいて重要となる「具体的な形状（ジオメトリ）」や「空間的な構成」を伝える媒体としては、言語は本質的に不向きである。（例：「右側をなめらかに隆起させて」といった微細な形状操作を、テキストだけで正確に指示することは困難である）。
    \item この「伝えたい形状イメージ（非言語情報）」と「入力手段（言語）」の間のギャップが、意図通りの生成を妨げる最大の要因となっている。
  \end{itemize}

  \item \textbf{「こだわり」の反映手段の欠如}
  \begin{itemize}
    \item 既存の Image-to-Image（画像検索による入力）等の手法もあるが、ユーザーが頭の中に描いている独自の形状特徴（こだわり）と完全に一致する画像が、オンライン上に存在するとは限らない。
    \item 結果として、ユーザーは形状に対する具体的な意図を持っているにも関わらず、それをAIに伝達する手段を持たないため、多義的な言葉（「かっこいい」等）に頼らざるを得ず、偶然の正解を待つ非効率な試行錯誤を強いられている。
  \end{itemize}

  \item \textbf{求められるアプローチ}
  \begin{itemize}
    \item 以上の分析から、生成AIを用いたデザイン探索においては、言語による意味的な制御に加え、ユーザーの意図する形状特徴を直接的にシステムへ伝達するための、非言語的な入力モダリティの統合が必要不可欠であると言える。
  \end{itemize}
\end{itemize}


\subsection{探索的デザイン支援システム}
\begin{itemize}
  \item 本研究では物理的な形状入力と対話的な探索を融合したデザイン支援システムを提案する。
  \item 本システムは、粘土による形状入力、LLMによる多義性解消（解釈提案）、および制約に基づく画像生成探索の3段階で構成される。
  \item これにより、形状と言語の両面からユーザーのメンタルモデルに寄り添ったデザイン生成を実現する。
  \item このために、以下のハイブリッドなアプローチをとる。
  \begin{itemize}
    \item 1.物理的インターフェース（粘土）による非言語的な意図伝達: 言語化できない形状のニュアンス（曲がり具合、ボリューム感など）を、粘土という可塑性のある物理メディアを通じて直接入力する。これにより、言語プロンプトの限界を補完する。
    \item 2.構造的解釈による多義性の解消: ユーザーの入力（画像・テキスト）に含まれる曖昧さを、AIが「モチーフ（何を作るか）」と「属性（どう作るか）」に構造化して解釈する。システムが複数の解釈案を提示し、ユーザーがそれを選択することで、デザインの方向性を合意形成する。これにより、ユーザーの不明瞭なイメージを言語化(外在化)し、AIとの解釈との乖離を低減する。
    \item 3.制約に基づく探索的ナビゲーション: 一発で正解を出そうとするのではなく、「もっと鋭く」「もっと有機的に」といった相対的な制約を与えることで、広大なデザイン空間の中からユーザーの好む領域を探索的に絞り込んでいく。
    % \item 言語化困難な「スタイルや詳細」を、AIからの提案（解釈）と選択（制約）によって外在化させるため。
  \end{itemize}
  \item この探索プロセスを実現するために、本研究では検索支援システムFindMeで提案された『制約に基づくナビゲーション』の概念を応用する。
  \item FindMeは既存のカタログからの検索を対象としていたが、本システムではこれを生成モデルの潜在空間探索へと拡張し、『もっと〇〇に』という比較制約による直感的なデザイン生成を実現する。
  \item また、意図した範囲内（制約内）で多様な候補を提示することで、ユーザー自身も気づいていなかった好みの発見（発想の拡張）を促す。
  \item 最終的に、非専門家ユーザーのメンタルモデルとの整合性が高いデザイン生成を支援することを目的とする。
  \item 先行研究であるStyleCLIPなどは、テキスト入力によって画像の特徴を操作する手法を提案している。しかし、これらは計算機的な整合性（テキストといかに一致するか）を重視しており、ユーザー個人の『メンタルモデルとの整合性』や『主観的な満足度』を直接支援するものではない。そのため、非専門家が自身の曖昧なイメージを具体化する用途には適していない。
\end{itemize}





\section{デザインプロセスの設計}
\begin{itemize}
  \item 本章では、提案手法を用いた具体的なデザインプロセスの設計について述べる。
  \item 本プロセスは、ユーザーの曖昧なイメージを段階的に明確化する「意図の詳細化」と、その意図の範囲内で最適解を探る「探索」の反復によって構成される。
\end{itemize}

\subsection{デザインプロセスの流れ}
\begin{itemize}
  \item 本システムにおけるデザイン探索は、大きく「意図の詳細化（Refinement）」と「空間内の探索（Exploration）」の2段階のフェーズを繰り返すことで進行する。
  \item ユーザーは編集対象の粒度に応じて「全体編集」「部分編集」「合成」のいずれかを選択し、以下のフローを経てデザインを行う。
  % \item ユーザーは最初に編集モード（全体編集・部分編集・合成）を選択し、以下のプロセスを経てデザイン探索を行う。
  % \item ユーザーの曖昧なクエリ入力を、画像生成AIが解釈可能な具体的な指示へと変換する。
  % \begin{itemize}
  % \item \textbf{解釈の生成}: システムは入力されたクエリや画像に対し、複数の解釈案（形状の変化、接合方法のニュアンス等）を提示する。
  % % \item \textbf{モチーフ抽出}: クエリ内に固有名詞（例：「チューリップのように」）が含まれる場合、これを「モチーフ」として抽出し、属性情報とは分離して管理する。
  % \item \textbf{ユーザー選択}: 提示された解釈案から、ユーザーは自身の意図に最も近いものを選択する。満足する案がない場合は再生成を行う。
  % \end{itemize}
  % \item 選択された解釈に基づき、デザインを具体的な画像として可視化する。
  % \begin{itemize}
  % \item \textbf{特徴ベクトル変換}: 選択された解釈（自然言語）を、定義済みの属性空間上の点（特徴ベクトル）へと変換する。部分編集の場合は、当該パーツに対応する部分ベクトルを生成する。
  % \item \textbf{画像生成}: 特徴ベクトルおよび抽出されたモチーフ情報をプロンプトに組み込み、画像生成エンジンを用いてデザイン案を提示する。
  % \end{itemize}
  % \item 各モードに応じた初期入力を行い、システムは入力を解析可能な形式に構造化する。
  % \item \textbf{制約の付与}:生成された画像を見て、ユーザーは特定の属性に関する「制約」（例：「もっと鋭利に」）を追加し、特徴ベクトルの値を更新し画像を再生成することを繰り返すことでデザインを洗練させる。
  % \item 生成画像を元に手元の粘土模型こねて再び入力する旨も書く
  \item \textbf{1. マルチモーダル入力と構造化（Phase A）}:
  ユーザーは粘土模型の画像と、その形状に対する意図（テキスト）を入力する。システムは選択されたモード（全体・部分・合成）に応じて入力を解析可能な形式に構造化する。

  \item \textbf{2. 意図の解釈と合意形成（Phase B）}:
  システムは入力された多義的なクエリに対し、形状や接合方法に関する複数の「解釈案」を提示する。ユーザーはその中から自身のメンタルモデルに最も近いものを選択することで、デザインの方向性を決定する（意図の詳細化）。

  \item \textbf{3. 特徴ベクトルへの変換と生成（Phase C）}:
  選択された解釈に基づき、システムは自然言語の指示を「属性空間」上の点（特徴ベクトル）へと変換し、具体的なデザイン案として画像を生成・提示する。

  \item \textbf{4. 制約に基づく探索的洗練（Phase D）}:
  生成された画像に対し、ユーザーは「もっと鋭利に」といった属性レベルの「制約」を追加する。システムはこの制約を満たす範囲内で特徴ベクトルを更新し、画像を再生成する。このサイクルを繰り返すことで、細部のニュアンスを探索する。

  \item \textbf{5. 物理空間へのフィードバック}:
  デジタル上で満足のいくデザインが得られた場合、ユーザーはその画像を参考に手元の粘土模型を修正し、再度システムに入力することで、より高精度なデザイン探索へと移行する。
\end{itemize}

\subsection{デザイン意図の詳細化}
% \begin{itemize}
%   \item フェーズA: マルチモーダル入力と構造化
%   \item モチーフ抽出: クエリ内に固有名詞（例：「チューリップのように」）が含まれる場合、これを「モチーフ」として抽出し、属性情報(後述)とは分離して管理する。
%   \item 理由として「チューリップ」のように」は具体的な形状イメージを指している可能性が高く、属性空間で表現される抽象的な特徴とは異なるため。
%   \item フェーズB: 意図の解釈と提案 (Interpretation)
%   \item まずユーザーが作成した粘土のラフモデルを画像生成AIへの形状の指示情報として利用する。
%   \item その後抽象的なクエリが持つ多義性を解消するために、「解釈の提案」を行う。
%   \item 非専門家の入力するクエリ（例：「和風」）は、多様な解釈が可能である。
%   \item システムは、そのクエリに関連する代表的な視覚的特徴（属性の組み合わせ）を複数パターン生成し、ユーザーに提示する。
%   \item 例：解釈A：[属性：茶色、低彩度、木質]（落ち着いた和風）, 解釈B：[属性：赤、白、金]（祝祭的な和風）
%   \item 属性：「金属的」「角が丸い」といった具体的な特徴
%   \item 属性値：その適合度合いを表す0 から1 の実数値。
%   \item ユーザーがいずれかの解釈を選択することで、システムは予め設計した数グループの属性空間を基にユーザーが求めている「和風」の具体的なパラメータ（特徴ベクトル）を特定し、方向性の合意形成を行う。
% \end{itemize}
\begin{itemize}
  \item \textbf{フェーズA: マルチモーダル入力と構造化}
  \begin{itemize}
    \item デザインの初期段階において、ユーザーの持つイメージは言語化が困難で曖昧である。
    \item 本節では、この曖昧な入力をシステムが処理可能な「明確な要求仕様」へと変換するプロセスについて述べる。
    \item ユーザーは粘土画像に加え、テキストによる指示を入力する。この際、システムは以下の3つのモードに応じた構造化を行う。
    \begin{itemize}
      \item \textbf{全体編集}: 画像全体を入力とし、物体全体の形状特徴を定義する。
      \item \textbf{部分編集}: 画像に加え、編集したい箇所をマスク画像で指定し、その部位に対する局所的な変更指示を受け付ける。
      \item \textbf{合成}: 2つのパーツ（部品と本体）が写った画像を入力し、それらの接合方法や配置に関する指示を受け付ける。
    \end{itemize}
    \item \textbf{モチーフ抽出}: クエリ内に固有名詞（例：「チューリップのように」）が含まれる場合、これを「モチーフ」として抽出し、後述する属性情報とは分離して管理する。これは、固有名詞が具体的な形状イメージ（What）を指している可能性が高く、属性空間で表現される抽象的な様式特徴（How）とは区別して扱う必要があるためである。
  \end{itemize}

  \item \textbf{フェーズB: 意図の解釈と提案 (Interpretation)}
  \begin{itemize}
    \item 非専門家が入力する抽象的なクエリ（例：「和風にして」）は多義的であり、一意な形状操作に変換することが困難である。そこでシステムは「解釈の提案」を行うことで多義性を解消する。
    \item システムは粘土画像を視覚的コンテキストとして利用しつつ、クエリに関連する代表的な属性の組み合わせ（解釈）を複数パターン生成し、ユーザーに提示する。
    \item (合成モードの場合は、クエリから接合方法に関する解釈を生成し提示する)
    \item \textbf{提示例}:
    \begin{itemize}
      \item 解釈A（落ち着いた和風）: [属性: 茶色, 低彩度, 木質]
      \item 解釈B（祝祭的な和風）: [属性: 赤, 白, 金, 装飾的]
    \end{itemize}
    \item ここで「属性」とは「金属的」「角が丸い」といった具体的な視覚的特徴を指し、「属性値」はその適合度合いを0から1の実数値で表したものである。
    \item ユーザーがいずれかの解釈を選択することで、システムは予め設計された属性空間の中からユーザーが求めている具体的なパラメータ（特徴ベクトル）を特定し、デザインの方向性に関する合意形成を行う。
  \end{itemize}
\end{itemize}

\subsection{デザイン意図に基づく探索}
\begin{itemize}
  \item 意図の方向性が定まった後も、細部の調整においては言語化が難しい微細な意図が存在する。
  \item 本節では、属性空間内での局所的な探索によってデザインを洗練させるプロセスについて述べる。
  \item \textbf{フェーズC: 特徴ベクトル生成と画像生成}
  \begin{itemize}
    \item フェーズBで合意された解釈（自然言語）を、画像生成エンジンが制御可能な数値表現である「特徴ベクトル」へと変換する。
    \item この際、編集モードに応じて以下の2種類のベクトルを使い分ける。
    \begin{itemize}
      \item \textbf{グローバルベクトル}: 画像全体の特徴（全体編集・合成で使用）
      \item \textbf{部分ベクトル}: 特定の部位の特徴（部分編集で使用）。マスクされた領域に対してのみ適用される。
    \end{itemize}
    \item 特徴ベクトルおよび抽出されたモチーフ情報をプロンプトに組み込み、画像生成エンジンを用いて具体的なデザイン案を提示する。
    \item (合成モードの場合、画像生成エンジンを用いて具体的なデザイン案を提示し、そのデザイン案に基づいて特徴ベクトルを生成する)
  \end{itemize}

  \item \textbf{フェーズD: 制約探索によるデザイン洗練}
  \begin{itemize}
    \item 生成された画像を見ながら、ユーザーは特定の属性に対して「制約」を追加することで、探索空間を絞り込む。
    \item \textbf{制約の付与}: 例えば「もっと鋭利にしたい」という場合、ユーザーは「鋭利さ（Sharpness）」という属性に対し、「$P \ge 0.8$」のような不等号を用いた制約条件を追加する。
    \item \textbf{空間内の再探索}: システムは、現在の特徴ベクトル周辺において、追加された制約条件を満たす部分空間を特定し、その中でランダムに特徴ベクトルを更新して画像を再生成する。
    \item これにより、意図を満たしつつも多様なバリエーション(デザイン案)を提示することで、ユーザー自身も想定していなかった好みの発見や、発想の拡張を促す。
    % \item これにより、ユーザーは言語化しにくい「程よいバランス」を探索的に発見できる。
    \item \textbf{斥力による多様性維持}: 特徴ベクトルを更新する際、過去にユーザーが「好ましくない」として棄却（クローズ）したデザイン案のベクトルからの「斥力」を作用させる。
    \item 斥力は新しく展開する特徴ベクトルが、過去の失敗作に近づかないようにする力として機能する。
    \item 具体的には、過去にクローズされたノードのベクトルとの差分ベクトルを計算し、その逆方向に新しいベクトルをシフトさせることで実現する。
    \item これにより、似通った失敗作の再生成を防ぎ、常に新しい可能性を提示することで探索の多様性を確保する。
  \end{itemize}
\end{itemize}
% \begin{itemize}
%   \item フェーズC: 特徴ベクトル生成と画像生成
%   \item フェーズD: 制約探索によるデザイン洗練
%   \item 次に、「制約」を用いた探索によって、形状や細部のニュアンスを調整する。
%   \item 方向性が定まった後も、細部の調整においては言語化が難しい微細な意図が存在する。
%   \item 生成された画像を見ながら、ユーザーが特定の属性に対して「制約」を追加する。
%   \item 制約：属性値P$\le$0.6のように不等号を用いて属性の適合度合いを制限する
%   \item システムはユーザーが追加した制約を満たす空間内でランダムに特徴ベクトルを更新
%   \item 多様性維持: 更新の際、過去にユーザーが棄却（クローズ）したノードからの「斥力」を作用させることで、似通った失敗作の再生成を防ぎ、探索の多様性を確保する。
%   \item 探索の際は、2種の特徴ベクトルを用いる
%   \begin{itemize}
%     \item グローバルベクトル: 画像全体の特徴を表すベクトル
%     \item 部分ベクトル: 特定のパーツ/部位の特徴を表すベクトル。特定のパーツの調整を行う際にマスク作成+クエリ入力→クエリ解釈を元に生成し、制約により更新する。
%   \end{itemize}
% \end{itemize}



\subsection{詳細化と探索の繰り返しによるデザイン}
\begin{itemize}
  \item \textbf{デジタル空間での探索ループ}
  \begin{itemize}
    \item ユーザーは生成画像に対し、「全体的な印象の変更（フェーズBへの戻り）」や「細部の微調整（フェーズDのループ）」、あるいは「特定部位の修正（部分編集モードへの切り替え）」を繰り返すことで、メンタルモデルとの整合性を高めていく。
  \end{itemize}

  \item \textbf{物理とデジタルの相互作用}
  \begin{itemize}
    \item 本システムはデジタル上の完結を目的とせず、デジタル上で意図に合致した、あるいは新たな着想の源となる生成画像が得られた段階で、ユーザーはその視覚情報を参考に手元の粘土模型を物理的に修正する。
    \item 修正された粘土模型を再度「フェーズA」に入力することで、システムは更新された形状をベースラインとして認識し、より解像度の高いデザイン提案が可能となる。
    \item この「物理入力 $\to$ AIによる拡張 $\to$ 物理修正 $\to$ 再入力」というサイクルの繰り返しにより、ユーザーの曖昧な初期イメージは、実装可能なレベルの具体的なデザインへと昇華される。
  \end{itemize}
\end{itemize}
% \begin{itemize}
%   \item ＊詳細化と探索を繰り返すことなどを書く＊
%   \item 生成された画像を起点として、ユーザーはさらなる修正や微調整を行う（探索ループ）。
%   \item 生成画像を元に手元の粘土模型こねて再び入力する旨を書くのかな？
% \end{itemize}


\begin{figure}[H]
  \centering
  \includegraphics[width=0.65\linewidth,keepaspectratio]{/Users/umeking5/university/lab_thesis/flows/flow_entire.png}
  \caption{本システムの全体フロー}
  \label{fig:system-flow}
\end{figure}

\begin{figure}[H]
  \centering
  \includegraphics[width=\textwidth]{/Users/umeking5/university/lab_thesis/flows/flow_a.drawio.png}
  \caption{フェーズA：入力}
  \label{fig:system-flow}
\end{figure}

\begin{figure}[H]
  \centering
  \includegraphics[width=\textwidth]{/Users/umeking5/university/lab_thesis/flows/flow_bc.drawio.png}
  \caption{フェーズB：解釈提案 / フェーズC：特徴ベクトル生成と画像生成}
  \label{fig:system-flow}
\end{figure}

\begin{figure}[H]
  \centering
  \includegraphics[width=\textwidth]{/Users/umeking5/university/lab_thesis/flows/flow_d.drawio.png}
  \caption{フェーズD：制約探索によるデザイン洗練}
  \label{fig:system-flow}
\end{figure}






\section{実験例と考察}
\begin{itemize}
  \item 本章では、提案手法（粘土による形状入力と生成AIを活用した探索）が、デザインの専門知識を持たないユーザーにとって、メンタルモデルとの整合性が高い画像を生成するうえで有効であるかを検証するための実験について述べる。
  \item 具体的には、従来のプロンプト入力のみの手法と比較し、探索の効率性（定量的評価）および、ユーザーの主観的な満足度や発想支援効果（定性的評価）が向上することを確認することを目的とする。
\end{itemize}

\subsection{システム構成}
\begin{itemize}
  % \item 本システムはデジタル上の完結を目的とせず、ユーザーが得られた生成画像を参考に手元の粘土模型を物理的に修正する。
  % \item 修正後の模型を再度システムに入力することで、デジタルとフィジカルを行き来する反復的なデザイン支援を実現する。
  \item 本システムは、以下の技術スタックおよびデータ構造に基づいて構築されている。
  \item \textbf{実装技術}: Streamlit (UI), OpenAI GPT-4o (推論・解釈), GPT-Image-1 (画像生成), Inpainting API (部分編集・合成)
  \item \textbf{属性空間}: 形状や質感を表す言語的な特徴量（Attribute Vector）を定義した多次元空間。
\end{itemize}

\subsection{実験方法}
\begin{itemize}
  \item 属性空間：形状に関する属性53属性
  \item タスク設定：正解画像を用意せず抽象的なお題を3題用意する。非専門家の利用を考慮し、日用品を対象とする。
  \item 題1「椅子をデザインしてください」
  \item 題2「カップをデザインしてください」
  \item 各条件でユーザーが満足するまで生成、時間無制限
  \item 提案手法: 形状入力とテキストによる初期入力。解釈提案と特徴ベクトル生成を経て、制約による微調整を行いながら反復的に画像生成
  \item 比較対象設定1: 形状入力とテキストによる初期入力。テキストプロンプトのみで画像反復生成。反復して画像を生成する際直前の生成デザインを入力。
  \begin{itemize}
  \item 提案手法と比較し、「特徴ベクトル」や「制約」がないと微調整で全体が崩れてしまい、制御性が低いことを証明
  \end{itemize}
  \item 比較対象設定2: テキストによる初期入力のみ。テキストプロンプトのみで画像反復生成。反復して画像を生成する際直前の生成デザインを入力しない。
  \begin{itemize}
  \item 比較対象1と比較し、て形状入力がないとユーザーの意図する形状特徴を反映しにくいことを証明
  \end{itemize}
  \item 定量的評価：座標の推移,収束までに生成した画像枚数などを計測し、効率性を比較する
  \item 定性的評価：作業中や作業後のインタビュー、作業中の発言内容の分析により、システムの使用感、デザインの満足度を測る
\end{itemize}
\subsection{結果の説明}
\subsection{考察}




\section{結論}
\begin{itemize}
  \item 形状入力と、生成AIを用いたクエリの解釈提案と制約を用いた探索空間の制限によって非専門的なユーザーでもメンタルモデルとの整合性の高いデザインの創作が支援され、簡単になることを目指した
  \item 結果考察から得られるシステムの改良点や今後の展望について述べる
\end{itemize}



\end{document}
